\documentclass[10pt,a4paper]{article}
\usepackage[margin=19mm]{geometry}
\usepackage{amsmath,amssymb,booktabs,tabularx,array,xcolor,placeins,hyperref}
\hypersetup{colorlinks=true,linkcolor=blue!50!black,pdftitle={Table II(a): Resident--streaming demand at two residency endpoints}}
\setlength{\parindent}{0pt}
\setlength{\parskip}{7pt}
\setlength{\emergencystretch}{3em}
\renewcommand{\thetable}{II(a)}
\title{\bfseries Residency Endpoints for Logical Matrix--Vector Tasks}
\author{}
\date{24 September 2026}
\begin{document}
\maketitle
\vspace{-1em}
\textbf{Reference.} A logical weight matrix $W\in\mathbb{Z}^{N\times K}$ is partitioned over resident $128\times128$ service units. An input slice is counted once at each receiving tile, so one vector evaluation receives
\[
Q_S=\left\lceil\frac{N}{128}\right\rceil K\ \mathrm{B}.
\]
Each valid resident weight element is written once when loading is included:
$Q_R=NK\ \mathrm{B}$ for per-use reloaded, and $Q_R=0$ for weight-static.

% Generated by shared/scripts/generate_IIa.py; do not hand-edit values.
% Requires booktabs, tabularx, array. Standalone numbering is set by table_IIa.tex.
\begin{table*}[htbp]
\centering
\caption{Resident--streaming demand at two residency endpoints.}
\label{tab:taskIIa-residency}
\small
\setlength{\tabcolsep}{3pt}
\renewcommand{\arraystretch}{1.2}
\begin{tabularx}{\textwidth}{@{}l*{5}{>{\centering\arraybackslash}X}@{}}
\toprule
Residency & $128\!\times\!128$ & $1024\!\times\!1024$ & $4096\!\times\!4096$ & $1024\!\times\!4096$ & $4096\!\times\!1024$ \\
\midrule
Weight-static & \begin{tabular}[t]{@{}c@{}}$Q_S=128\,\mathrm{B}$\\$Q_R=0\,\mathrm{B}$\\$\mathrm{RI}=\infty$\end{tabular} & \begin{tabular}[t]{@{}c@{}}$Q_S=8\,\mathrm{KiB}$\\$Q_R=0\,\mathrm{B}$\\$\mathrm{RI}=\infty$\end{tabular} & \begin{tabular}[t]{@{}c@{}}$Q_S=128\,\mathrm{KiB}$\\$Q_R=0\,\mathrm{B}$\\$\mathrm{RI}=\infty$\end{tabular} & \begin{tabular}[t]{@{}c@{}}$Q_S=32\,\mathrm{KiB}$\\$Q_R=0\,\mathrm{B}$\\$\mathrm{RI}=\infty$\end{tabular} & \begin{tabular}[t]{@{}c@{}}$Q_S=32\,\mathrm{KiB}$\\$Q_R=0\,\mathrm{B}$\\$\mathrm{RI}=\infty$\end{tabular} \\[5pt]
\shortstack[l]{Per-use\\reloaded} & \begin{tabular}[t]{@{}c@{}}$Q_S=128\,\mathrm{B}$\\$Q_R=16\,\mathrm{KiB}$\\$\mathrm{RI}=1/128$\end{tabular} & \begin{tabular}[t]{@{}c@{}}$Q_S=8\,\mathrm{KiB}$\\$Q_R=1\,\mathrm{MiB}$\\$\mathrm{RI}=1/128$\end{tabular} & \begin{tabular}[t]{@{}c@{}}$Q_S=128\,\mathrm{KiB}$\\$Q_R=16\,\mathrm{MiB}$\\$\mathrm{RI}=1/128$\end{tabular} & \begin{tabular}[t]{@{}c@{}}$Q_S=32\,\mathrm{KiB}$\\$Q_R=4\,\mathrm{MiB}$\\$\mathrm{RI}=1/128$\end{tabular} & \begin{tabular}[t]{@{}c@{}}$Q_S=32\,\mathrm{KiB}$\\$Q_R=4\,\mathrm{MiB}$\\$\mathrm{RI}=1/128$\end{tabular} \\[5pt]
\bottomrule
\end{tabularx}
\par\smallskip
\begin{minipage}{\textwidth}\footnotesize
The primary boundary sums valid logical payload at all $128\times128$ tile input receivers. Each column specifies the logical weight shape $N\times K$ (outputs $\times$ inputs), not a new physical macro. Each window evaluates one input vector; input and resident elements are one byte (INT8 reference). Weight-static assumes preloading before the window; per-use reloaded writes the complete weight matrix once within it. $\mathrm{RI}=Q_S/Q_R$, with $\mathrm{RI}=\infty$ for $Q_R=0$. Here, all dynamic ratios equal $1/128=0.0078125$. $1\,\mathrm{KiB}=1024\,\mathrm{B}$; $1\,\mathrm{MiB}=1048576\,\mathrm{B}$.
\end{minipage}
\end{table*}

\FloatBarrier

\textbf{Interpretation.} All five shapes have output dimensions divisible by 128. Their local-receiver ratios therefore coincide at $1/128$ despite different demand volumes and tile counts. The two rectangular tasks also have the same local-receiver $Q_S$ and $Q_R$ in this reference.

The contrast boundary, documented in the Chinese method and machine data, counts the whole-matrix input only once: $Q_S^{\mathrm{op}}=K\ \mathrm{B}$ and $\mathrm{RI}^{\mathrm{op}}=1/N$ for per-use reloaded. It is not mixed into the table cells above.

\textbf{Scope.} One use is one vector evaluation. Per-use reloaded is neither an entire training step nor one load amortized over a multi-vector GEMM. Summed receiver demand is not a runtime prediction obtained by dividing by a single macro's throughput. A future hardware comparison must match the service-unit set, write mode and precision.
\end{document}
